% Глава: 11-mc
% Статус: наполнено (первый прогон).
\chapter{Файловый менеджер: Midnight Commander}
\label{ch:11-mc}

\section{Зачем нужен}

\cmd{mc} (Midnight Commander) --- терминальный файловый
менеджер с «двухпанельным» интерфейсом. Полезен, когда
файловые операции проще сделать «глазами», чем командами:
сравнить две директории, переместить группу файлов,
просмотреть файл встроенным просмотрщиком.

Полезен, когда вы работаете с файлами на своей ВМ на
\fpath{pglab.nsu.ru} и нужен быстрый «визуальный» контроль.

\begin{figure}[htb]
  \centering
  \begin{tikzpicture}[node distance=1.4cm]
    \node[box, align=left, minimum width=4.5cm, minimum height=3.5cm] (left)
      {ЛЕВАЯ ПАНЕЛЬ\\
       \texttt{incoming/}\\
       \texttt{a.txt}\\
       \texttt{b.txt}};
    \node[box, align=left, right=2.5cm of left, minimum width=4.5cm,
          minimum height=3.5cm] (right)
      {ПРАВАЯ ПАНЕЛЬ\\
       \texttt{processed/}\\
       \texttt{a.txt}};
    \node[below=1.2cm of left, font=\small] {командная строка и меню внизу};
  \end{tikzpicture}
  \caption{Схема интерфейса Midnight Commander.}
  \label{fig:mc}
\end{figure}

\section{Навигация}

Внутри \cmd{mc}:
\begin{itemize}
  \item \cmd{Tab} --- переключить активную панель;
  \item стрелки --- перемещение по списку;
  \item \cmd{Enter} --- войти в каталог;
  \item \cmd{Backspace} / \cmd{Alt+u} --- вверх;
  \item \cmd{Ctrl+t} --- создать каталог;
  \item \cmd{Alt+s} --- показать свойства файла;
  \item \cmd{Alt+g} --- перейти по пути вручную;
  \item \cmd{Alt+z} / \cmd{Alt+x} --- сжать / распаковать архив;
  \item \cmd{F3} --- просмотр файла, \cmd{F4} --- открыть в редакторе.
\end{itemize}

\section{Файловые операции}

Выделяйте файлы клавишей \cmd{Insert} (или \cmd{Space}) и
используйте \cmd{F5} (копировать), \cmd{F6} (переместить/
переименовать), \cmd{F8} (удалить). Все операции
выполняются с подтверждением и показывают прогресс.

Полезные сочетания:
\begin{itemize}
  \item \cmd{F7} --- создать каталог;
  \item \cmd{F8} --- удалить;
  \item \cmd{F6} --- переместить или переименовать;
  \item \cmd{F5} --- копировать;
  \item \cmd{Ctrl+r} --- сравнить две директории (включая
    содержимое).
\end{itemize}

\section{Поиск и сортировка}

\cmd{Alt+s} --- свойства, \cmd{Ctrl+s} / \cmd{Alt+s} ---
поиск по имени; \cmd{Alt+s} --- переключить сортировку
(по имени, дате, размеру). \cmd{Ctrl+h} --- показать/
скрыть скрытые файлы (начинающиеся с \cmd{.}).

\section{Когда mc, а когда shell-команды}

\begin{center}
\begin{tabular}{lp{0.55\linewidth}}
\toprule
\textbf{Задача} & \textbf{Удобнее} \\
\midrule
одиночный файл & \cmd{cp} / \cmd{mv} \\
копировать дерево & \cmd{cp -r} \\
просмотр дерева & \cmd{mc} / \cmd{tree} \\
фильтр по имени & \cmd{find} + \cmd{grep} \\
сравнить две директории & \cmd{mc} (Ctrl+r) \\
пакетная обработка & shell-скрипт \\
\bottomrule
\end{tabular}
\end{center}

\cmd{mc} не заменяет shell-команды --- это другой интерфейс
к тем же операциям. Для повторяющихся задач скрипт на
оболочке (\nameref{ch:27-scripting-ide}) масштабируете,
а \cmd{mc} --- для разовых манипуляций.

\section{Типичные ошибки}
\begin{itemize}
  \item Не видите скрытые файлы --- включите
    \cmd{Ctrl+h} (или настройку «Show hidden files»).
  \item Переместили файл не туда --- посмотрите, какая
    панель была \emph{активной} (обведённая рамкой);
    операция выполняется из активной в неактивную.
  \item Не выходите из \cmd{mc}, нажав \cmd{q} ---
    используйте \cmd{F10} или \cmd{Ctrl+q}.
\end{itemize}

\section{Связь с остальным курсом}

\cmd{mc} удобен для навигации по структуре, созданной
пакетами (\nameref{ch:08-packages}) или Docker
(\nameref{ch:21-docker-data}): в контейнере, где нет
\cmd{tree}, \cmd{mc} часто единственный графический
по терминалу интерфейс для файлов.

\section{Задачи для самопроверки}
\begin{enumerate}
  \item Как переключаться между левой и правой панелью?
  \item Как создать в активной панели каталог
    и перейти в него?
  \item Как скопировать несколько выделенных файлов
    в другую директорию?
  \item Как сравнить две директории на различия?
  \item Чем работа в \cmd{mc} отличается от \cmd{find}
    + \cmd{grep} по гибкости?
\end{enumerate}
